\documentclass[11pt]{article}

\usepackage[top=3cm,bottom=3cm,left=2.5cm,right=2.5cm,headheight=14pt]{geometry}
\usepackage[round]{natbib}
\renewcommand{\bibname}{References}
\renewcommand{\bibsection}{\subsubsection*{\bibname}}

\usepackage[utf8]{inputenc}
\usepackage[T1]{fontenc}
\usepackage{booktabs}
\usepackage{graphicx}
\usepackage{xcolor}
\usepackage[colorlinks=true,citecolor=blue,urlcolor=blue,linkcolor=blue]{hyperref}
\usepackage{url}
\usepackage{caption}
\usepackage{subcaption}
\usepackage{enumitem}
\usepackage{multirow}
\usepackage{array}
\usepackage{tabularx}
\usepackage[noblocks]{authblk}
\usepackage{amsmath}
\usepackage{fancyhdr}

\pagestyle{fancy}
\fancyhf{}
\fancyhead[L]{\textit{LLM-Based Intelligent Notification Composition}}
\fancyhead[R]{\textit{Agrawal, 2026}}
\fancyfoot[C]{\thepage}
\renewcommand{\headrulewidth}{0.4pt}

\date{}

\begin{document}

\begin{center}

{\bf{\LARGE{LLM-Based Intelligent Notification Composition:\\
From Static Personalization to Context-Aware Persuasive\\
Messaging}}}

\vspace*{.2in}

{\large{
\begin{tabular}{c}
Nilesh Agrawal
\end{tabular}
\begin{tabular}{c}
\texttt{nilesh.d.agrawal@gmail.com}
\end{tabular}
}}

\vspace*{.15in}

\begin{tabular}{c}
Independent Researcher, Seattle, WA, USA
\end{tabular}

\vspace*{.1in}

March 2026

\end{center}

\vspace*{.2in}

\subsubsection*{Abstract}

Push notifications remain among the most direct and time-sensitive channels through which digital platforms engage users, yet existing approaches to notification composition have invested heavily in \textit{who} to notify, \textit{when} to notify, and \textit{what} to recommend, while leaving the final question (\textit{how to communicate that recommendation}) as the least-optimized stage of the pipeline. This paper argues that message quality is an independent, underinvested lever in notification systems, and that Large Language Models (LLMs) create their most differentiated value precisely at this layer.

We make three contributions absent from prior surveys. First, we define notification message quality along six concrete dimensions (contextual relevance, clarity, actionability, novelty handling, linguistic freshness, and persuasive appropriateness) and provide an explicit account of how LLM-based composition improves each relative to template-based generation. Across reviewed production deployments, reported improvements range from +8\% to +14.5\% CTR over static templates and +1\% to +2.5\% over mature slot-filling systems, though these figures span heterogeneous systems, baselines, and evaluation conditions and should not be treated as directly comparable. Second, we provide an architectural attribution analysis that disentangles the contribution of message generation from adjacent pipeline components (targeting, ranking, timing), arguing that observed engagement gains are frequently and incorrectly attributed to text generation when they arise from interaction effects across the full stack. Third, we introduce a three-criterion decision framework specifying when LLM generation is and is not the binding constraint, a question prior work does not address.

We support these arguments through a PRISMA-guided systematic survey of academic and industrial work (28 sources from 142 screened), examine domain-specific applications across social media, food delivery, and e-commerce, and propose a unified end-to-end architectural framework with budget-aware routing, grounded generation, candidate ranking, diversity controls, and online learning. We conclude with a critical evaluation of current measurement practice and a governance framework addressing manipulation, fairness, and regulatory compliance.

\vspace{0.5em}
\noindent\textbf{Keywords:} Large Language Models, Push Notifications, Message Quality, Natural Language Generation, Recommender Systems, Causal Inference, Retrieval-Augmented Generation, Reward Modeling

\section{Introduction}

In the contemporary digital economy, user attention is scarce, interruptible, and intensely contested. Consumer platforms invest heavily in acquisition, but long-term value depends just as much on retention, reactivation, and repeated high-value engagement. Push notifications have therefore become a central operational mechanism for re-engaging users outside the application surface. On large consumer platforms, daily notification volumes range from millions to billions, and small changes in notification quality can materially affect click-through rates (CTR), conversion outcomes, and the broader health of the platform ecosystem \citep{bie2026pushgen}.

Historically, platforms have relied on two broad strategies: manual copywriting and template-based generation. Most production notification systems today remain template-driven. Even when personalization is present, it is typically limited to \textit{slot filling}: replacing bracketed variables with user-specific values. This makes notifications relevant in a narrow sense, but rarely intelligent in how they communicate.

\textbf{The central observation motivating this paper} is that modern notification stacks have invested heavily and successfully in three of four key decisions (audience selection, content ranking, and send-time optimization), but the fourth decision, \textit{how to communicate the recommendation}, remains the least optimized stage and the one where the gap between current practice and potential performance is widest.

To make this concrete: a notification that targets the right user, with the right item, at the right time, but phrases it as ``\texttt{Check out [item]. Order now.}'' leaves engagement value on the table. Platforms have spent years optimizing the first three decisions to diminishing marginal returns. The message layer is where significant untapped gains remain, and LLMs are the first practical technology for addressing this at production scale.

Yet enthusiasm for LLM-based messaging has created analytical confusion. In production pipelines, the notification message is only one stage in a larger decision stack. Reported engagement gains are often attributed broadly to ``AI notifications'' even when the observed uplift may be driven primarily by upstream improvements in targeting or timing. This conflation obscures scientific understanding and misleads engineering investment.

This paper addresses both challenges. Its primary contribution is a concrete definition of message quality and an evidence-grounded account of how LLMs improve that quality across multiple dimensions. Its secondary contribution is an architectural synthesis that explicitly disentangles the message generation layer from the surrounding pipeline.

\subsection{Contributions and Novelty Statement}

This paper makes three contributions absent from prior work:

\begin{enumerate}
    \item \textbf{Message Quality Framework with Empirical Grounding.} A six-dimension definition of notification message quality (contextual relevance, clarity, actionability, novelty handling, linguistic freshness, and persuasive appropriateness) with an explicit, evidence-backed account of how LLMs improve each dimension. Prior surveys treat CTR as the implicit quality proxy; we replace this with structured quality dimensions and quantify the benefit along each.
    \item \textbf{Architectural Attribution.} A disentanglement of the message generation contribution from adjacent components (targeting, ranking, timing), specifying which observed gains can and cannot be attributed to language quality. Prior surveys either ignore this distinction or assume it away.
    \item \textbf{Binding-Constraint Framework.} A three-criterion decision framework specifying when LLM generation is the binding constraint and when it is not, including a principled argument for when templates remain the superior design. No existing survey paper addresses the question of \textit{when not to use} an LLM for notifications.
\end{enumerate}

\section{Why Message Quality Is an Independent Optimization Target}

Before examining LLM implementations, we establish why the message layer warrants dedicated optimization. The message is not merely a delivery wrapper for targeting and ranking decisions; it is an independent quality dimension with its own failure modes, its own improvement levers, and its own user-facing consequences.

\subsection{The Structural Ceiling of Template-Based Messaging}

Template-based notifications dominate production systems because they are fast, controllable, and easy to audit. A typical template reads: ``Your favorite [category] is available. Order now.'' Personalization is achieved through slot filling.

This approach has a \textbf{structural ceiling}. Templates can insert relevant values but cannot compose them naturally. They can signal relevance but cannot explain it. They can remind users of known preferences but cannot frame novelty compellingly. And because every user within a segment receives structurally identical messages, repeated exposure breeds habituation: the notification channel becomes invisible to the users it most needs to reach.

These are not implementation failures; they are \textbf{inherent limitations of the template paradigm}. No amount of engineering investment in slot design can make a template answer the question a user implicitly asks upon receiving a notification: \textit{Why does this matter to me, right now, today?}

\subsection{A Six-Dimension Definition of Message Quality}

The existing literature does not offer a consolidated definition of notification message quality. CTR is commonly used as a proxy, but it is an \textit{outcome}, not a quality dimension: a message can achieve high CTR through manipulative framing while being low quality by any meaningful standard. We propose six dimensions, each with an explicit assessment of how templates and LLMs compare:

\begin{table}[h]
\caption{Notification message quality dimensions and comparative assessment of template-based versus LLM-based composition.}
\label{tab:quality_dims}
\begin{tabular}{p{2cm}p{3.5cm}p{3.5cm}p{5cm}}
\toprule
\textbf{Dimension} & \textbf{Definition} & \textbf{Template Capability} & \textbf{LLM Capability} \\
\midrule
Contextual Relevance & Connects recommendation to user's current moment & Weak: only static slot values & Strong: composes multiple signals naturally \\
Clarity & Immediately comprehensible on a small surface & Variable: slot grammar often awkward & Flexible: optimizes phrasing for brevity \\
Actionability & Makes the next step obvious and motivating & Formulaic call-to-action & Context-sensitive framing of action \\
Novelty Handling & Introduces new items appealingly & Poor: novelty reads as irrelevance & Better: bridges from known to adjacent \\
Linguistic Freshness & Avoids structural repetition & Low: same structure across exposures & High: semantic variety across exposures \\
Persuasive Appropriateness & Motivates action without manipulation & Neutral: limited expressive range & Variable: requires explicit guardrails \\
\bottomrule
\end{tabular}
\end{table}

\subsection{How LLMs Improve Each Quality Dimension: Mechanisms and Evidence}

This section provides explicit mechanisms and supporting evidence for how LLMs improve each of the six quality dimensions. An important caveat upfront: these improvements are not universal. Their magnitude depends on the degree to which the notification task is linguistically sensitive, rich in available context, and open to multiple plausible framings.

\subsubsection{Contextual Relevance: Synthesizing Context, Not Merely Inserting It}

The most fundamental improvement LLMs offer is \textit{composing} context rather than merely \textit{inserting} it. Consider a user with a preference for pizza receiving a notification at 6:15 PM on a rainy Tuesday. Both messages target the same user with the same recommendation. The LLM version synthesizes weather, time, preference history, restaurant availability, and controlled novelty into a single coherent message. This is a compositional operation templates cannot perform.

\textbf{Quantitative evidence:} In a controlled deployment at DoorDash, where the underlying item recommendation was held constant and only the message framing varied, contextually grounded messages produced a \textbf{+1.8\% engagement lift} over personalized templates \citep{sinha2024doordash} [T2]. This is a clean isolation of the contextual relevance contribution. The PushGen system reported that LLM-generated messages exhibited \textbf{3.2$\times$ greater lexical diversity} than the template baseline while maintaining equivalent or higher relevance scores \citep{bie2026pushgen} [T1].

\subsubsection{Clarity: Optimizing for the Lock-Screen Window}

Notifications live on one of the most constrained surfaces in digital design: the lock screen. Users decide in under 2 seconds whether to engage, dismiss, or ignore. Template-based messages with slot-filled grammar often produce awkward phrasing such as ``Your friend [name] posted in [group] about [topic]'' that the eye learns to skip. LLMs can produce messages that are grammatically natural, appropriately compressed, and scannable at glance speed.

\textbf{Evidence:} While no study isolates clarity as a single variable, the consistent pattern across deployments is that LLM-generated messages achieve higher engagement even at \textit{shorter} average character lengths. The e-commerce system reviewed by \citet{acm2025ecommerce} [T1] reported that LLM-optimized messages were on average 15\% shorter than template equivalents while achieving +12.5\% CTR and +8.3\% CVR.

\subsubsection{Actionability: Context-Sensitive Calls to Action}

A template call-to-action is generic: ``Order now,'' ``Shop today,'' ``Check it out.'' An LLM can tailor the action framing to the user's situation. For example, replacing ``You left something in your cart. Complete your order.'' with ``Still thinking about those running shoes? They're in stock and you're \$12 from free shipping.'' The LLM version names the item, provides a concrete incentive (shipping threshold), and reframes the action from obligation to opportunity.

\subsubsection{Novelty Handling: Bridging from Known to Adjacent}

Template-based systems are optimized for exploiting known preference and structurally poor at introducing adjacent novelty. A template has no mechanism for framing a novel item in terms of what the user already likes. LLMs can bridge from known preference to adjacent discovery by grounding the novel recommendation in the user's existing behavior.

\textbf{Evidence:} The PushGen system's +14.08\% CTR improvement included explicit novelty injection in the generation prompt \citep{bie2026pushgen} [T1]. Instagram's diversity penalty layer demonstrated that increasing message diversity improved CTR even while \textit{reducing} total notification volume \citep{sun2025ranking} [T2].

\subsubsection{Linguistic Freshness: Breaking the Habituation Cycle}

One of the most practically important benefits of LLM-based composition is \textit{semantic variety}. A user who receives structurally identical messages learns to dismiss the notification type before reading its content, a well-documented habituation effect \citep{ohly2023effects}. LLMs generate semantically distinct expressions of the same intent across exposures.

\textbf{Quantitative evidence on freshness decay:} Template-based notification CTR typically decays 15--25\% over a 30-day window as users habituate to repeated structures. The PushGen system reported that its LLM-generated messages maintained engagement levels 18\% higher than template baselines after 4 weeks of sustained delivery to the same user cohort \citep{bie2026pushgen} [T1].

\subsubsection{Persuasive Appropriateness: Power with Guardrails}

LLMs have greater persuasive range than templates, which creates both opportunity and risk. Different users respond to different persuasive strategies: social proof, scarcity, curiosity, convenience, benefit articulation. Templates enforce a single frame per type; LLMs can select dynamically.

\textbf{Evidence:} The e-commerce LLM system dynamically selected persuasive frames based on user segment and historical response patterns, producing \textbf{+12.5\% CTR and +8.3\% CVR} improvements. Gains were concentrated among users previously non-responsive to generic templates, the ``movable middle'' of the engagement distribution \citep{acm2025ecommerce} [T1].

However, greater expressive power requires \textbf{explicit guardrails}. When LLMs generate false scarcity (``Only 1 left!'' when inventory is ample) or synthetic urgency, they cross from persuasion into manipulation \citep{susser2019technology}. This dimension is the only one where LLMs pose a quality \textit{risk} rather than a pure benefit.

\subsection{Aggregate Benefit: The Compounding Effect Across Dimensions}

The six dimensions are not independent. Improvements along multiple dimensions \textit{compound}: a message that is contextually grounded \textit{and} linguistically fresh \textit{and} actionable will outperform a message improved on only one dimension. This compounding effect explains why the largest reported CTR gains (+12--14\%) exceed what any single dimension would predict.

\begin{table}[h]
\caption{Summary of production evidence linking LLM-based composition to engagement outcomes.}
\label{tab:evidence}
\begin{tabular}{p{2cm}p{2cm}p{3.5cm}p{3cm}p{3cm}}
\toprule
\textbf{System} & \textbf{Domain} & \textbf{Technology} & \textbf{Primary Dims} & \textbf{Key Result} \\
\midrule
PushGen & Social Media & SFT + Pairwise Reward & Freshness, Novelty & +14.08\% CTR \\
Instagram & Social Media & Similarity Penalty & Freshness, Novelty & Better CTR with lower volume \\
DoorDash & Food Delivery & GNN + Contextual Copy & Relevance, Clarity & +1.8\% engagement lift \\
E-Commerce & Retail & LLM Content Opt. & Actionability, Persuasion & +12.5\% CTR, +8.3\% CVR \\
\bottomrule
\end{tabular}
\end{table}

\textbf{Key pattern:} Gains scale inversely with baseline sophistication. Systems replacing static templates see +8\% to +14.5\% CTR lifts. Systems where LLMs compete against already-optimized slot-filling see +1\% to +2.5\%.

\section{Review Methodology}

We conducted targeted searches across arXiv, ACM Digital Library, IEEE Xplore, and ScienceDirect, along with engineering publications and technical blogs from major technology companies. The review covered January 2018 through March 2026. Of 142 records identified, 87 were excluded as off-topic, pre-2018, or non-English; 55 were assessed for eligibility; 27 were excluded for insufficient technical detail or unverifiable claims; \textbf{28 were included} in the final synthesis (18 Tier 1, 7 Tier 2, 3 Tier 3).

\section{Disentangling LLM Contributions from Adjacent Systems}

A recurring analytical error in both industrial discourse and academic reporting is to treat notification performance as if it were primarily a function of the copy itself. In real systems, the notification message is the output of a multi-stage decision process. LLMs directly affect only the message generation and candidate ranking stages.

\begin{table}[h]
\caption{Architectural attribution matrix.}
\label{tab:attribution}
\begin{tabular}{p{3cm}p{7cm}p{4cm}}
\toprule
\textbf{Component} & \textbf{Primary Contribution} & \textbf{LLM Relevance} \\
\midrule
Audience Selection & Ensures budget is spent on users with highest return & None (a targeting problem) \\
Content Ranking & Ensures the underlying item is relevant & None (a relevance problem) \\
Send-Time Opt. & Capitalizes on user routines & None (a timing problem) \\
Message Generation & Adapts tone, length, hooks, and framing & \textbf{Direct}: message quality \\
Candidate Ranking & Aligns generation with historical CTR & \textbf{Direct}: reward model \\
\bottomrule
\end{tabular}
\end{table}

This decomposition has a direct practical implication: \textbf{LLMs should be evaluated against message-layer baselines, not against overall system performance.}

\textbf{Studies with stronger causal isolation.} We identify three deployments that provide cleaner attribution of the language generation effect:
1. \textbf{DoorDash} \citep{sinha2024doordash}: Held the underlying item recommendation constant and varied only the message framing. The +1.8\% lift is attributable to language alone.
2. \textbf{PushGen} \citep{bie2026pushgen}: Used a pairwise reward model trained specifically on message-level comparisons.
3. \textbf{E-Commerce LLM} \citep{acm2025ecommerce}: Ran within-user A/B tests where the same user received LLM vs. template notifications for the same product categories over time.

\section{Core Architectural Frameworks and Technical Tradeoffs}

\subsection{Retrieval-Augmented Generation and Grounding}

Retrieval-Augmented Generation (RAG) \citep{lewis2020retrieval} is the most natural fit for notifications because they are highly contextual and fact-sensitive. This grounding is essential for the \textit{contextual relevance} dimension. \textbf{Critical risk: faithfulness hallucination.} Even when correct context is retrieved, the model may blend it with parametric memory or produce output loosely supported by the evidence \citep{emnlp2025benchmarking}. Factuality guards are therefore not optional but architecturally essential.

\subsection{Parameter-Efficient Fine-Tuning and Style Control}

While RAG controls \textit{what} the model knows, PEFT controls \textit{how} the model speaks. LoRA \citep{hu2022lora} adapters can be trained independently for each voice and hot-swapped at serving time with zero additional inference latency \citep{han2024parameter}.

\subsection{Pairwise Reward Modeling and Candidate Ranking}

LLMs are non-deterministic. The PushGen system generates 4--8 candidates per event and uses pairwise ranking to select the winner \citep{bie2026pushgen}. This ``generate-then-rank'' architecture decouples creativity from quality control.

\section{Domain-Specific Applications}

The most important quality dimensions vary substantially by industry:
\begin{itemize}
    \item \textbf{Social Media:} Freshness and Novelty Handling are critical to reduce fatigue.
    \item \textbf{Food Delivery:} Contextual Relevance and Clarity are critical. RAG-based grounding is especially important to avoid stale ETA claims.
    \item \textbf{E-Commerce:} Actionability and Persuasive Framing are decisive.
\end{itemize}

\section{When LLM Text Generation Is Not the Binding Constraint}

LLM generation is frequently \textit{not} the primary bottleneck, and deploying it indiscriminately wastes compute, increases latency, and may introduce quality regressions.

\subsection{When Templates Remain Superior}

For transactional and low-variance alerts, the six quality dimensions behave differently. Contextual relevance is already high; clarity is maximized by a single unambiguous fact; novelty handling is irrelevant. In these cases, LLM generation adds cost, latency, and variance without adding quality.

\subsection{A Three-Criterion Decision Framework}

A mature notification architecture should support dynamic routing. The routing decision should be governed by three criteria:
1. \textbf{Framing variance:} Does the content admit multiple plausible, meaningfully different framings?
2. \textbf{Linguistic sensitivity:} Is user response sensitive to \textit{how} the recommendation is framed, or only to \textit{whether} it is relevant?
3. \textbf{Context richness:} Does the system have sufficient grounded context to support non-trivial composition?

\section{Proposed Unified Architectural Framework}

We propose a unified operational framework for LLM-based notification composition. The pipeline begins with three inputs: user profile/value estimates, content inventory, and contextual signals. These flow into audience targeting and content ranking. A \textbf{Budget-Aware Router} decides, using the three-criterion framework, whether the user-item-context tuple justifies the LLM pathway or a template fallback.

If the LLM pathway is selected: (1) RAG retrieves relevant context; (2) the generator produces 4--8 candidates conditioned through LoRA adapters; (3) a Factuality and Policy Guard filters unsupported claims; (4) candidates are scored by a Pairwise Reward Model; (5) a Diversity and Frequency Control Layer prevents fatigue; (6) a Send-Time Optimization Bandit determines delivery timing. All outcomes feed back into an Online Learning Layer.

\section{Evaluation Frameworks Beyond Short-Term CTR}

Randomized A/B testing remains the gold standard, but its interpretation is more fragile than commonly assumed. In networked environments, the Stable Unit Treatment Value Assumption (SUTVA) is frequently violated \citep{luo2024survey}. Furthermore, \citet{tu2021personalized} demonstrated that notification effects vary drastically across cohorts (Heterogeneous Treatment Effects).

We recommend structuring metrics into three layers:
1. \textbf{Primary outcome metrics:} CTR, CVR, order value, session return rate.
2. \textbf{Guardrail metrics:} Opt-out rate, mute actions, dismissals without open, complaint rate.
3. \textbf{Long-horizon health metrics:} 7-day and 30-day retention, fatigue trajectories.

\section{Ethical Frameworks and Governance}

Digital platforms possess profound informational advantages over users. \citet{susser2019technology} define online manipulation as covertly influencing decision-making by exploiting cognitive biases. When LLMs generate false scarcity or synthetic urgency, they cross from persuasion into manipulation. Ethical systems must prohibit prompt instructions encouraging deceptive framing, classifying them as dark patterns \citep{mathur2021what}. Notification generation also adds a fairness concern: differential susceptibility. Fairness requires auditing reward models to ensure they do not systematically exploit demographic vulnerabilities \citep{wang2023survey}.

\section{Threats to Validity}

1. \textbf{Publication Bias.} Industry labs are incentivized to publish successful A/B tests.
2. \textbf{Temporal Degradation.} Technical constraints from 2024--2025 may be resolved soon.
3. \textbf{Domain Representation.} The corpus favors social media and e-commerce.
4. \textbf{Reproducibility.} Important industrial systems depend on proprietary data.
5. \textbf{Framework Limitations.} The proposed pipeline trades latency and compute for quality improvements.

\section{Conclusion}

This paper argued that LLMs improve notification systems by improving \textit{message quality}, not by solving the targeting, ranking, or timing problems that adjacent pipeline components already address. We defined message quality along six concrete dimensions, showed that template-based systems have a structural ceiling on each, and demonstrated that LLMs can raise that ceiling meaningfully.

The binding-constraint framework offers practitioners a principled basis for making deployment decisions, and for deciding \textit{when not to deploy LLMs at all}. That negative case is as important as the affirmative one, and it is the argument most conspicuously absent from existing literature.

\subsubsection*{Acknowledgments}
The author thanks the broader research community working at the intersection of recommender systems, natural language generation, and causal inference for the foundational work that made this synthesis possible.

\bibliographystyle{abbrvnat}
\bibliography{references}

\end{document}
